\documentclass[10pt,a4paper]{amsart}
\usepackage[margin=19mm]{geometry}
\usepackage{amsmath,amssymb,mathtools}
\usepackage[scr=rsfs,cal=euler]{mathalfa}
\usepackage{xfrac}
\usepackage[unicode,hidelinks,bookmarks=false]{hyperref}
\newcommand{\ie}{i.e.\ }
\newcommand{\cf}{cf.\ }
\newcommand{\resp}{respectively\ }
\newcommand{\Subsection}{\subsection}
\providecommand{\nobreakdash}{\nobreak\hbox{-}}
\setlength{\emergencystretch}{2em}
\multlinegap0pt
\usepackage{amssymb}
\usepackage{float}
\usepackage{mathtools}
\usepackage{enumerate}
\usepackage{xcolor}
\usepackage[shortlabels]{enumitem}
\usepackage{tikz}
\newtheorem{proposition}{Proposition}
\newtheorem{lemma}{Lemma}
\newtheorem{theorem}{Theorem}
\newtheorem{caselist}{Caselist}
\newcommand{\Lie}{\mathfrak{Lie}}
\newcommand{\Q}{\mathbb{Q}}
\newcommand{\QB}{\Q\langle B\rangle}
\newcommand{\ad}{\operatorname{ad}}
\newcommand{\lqq}{\mathfrak{lq}}
\newcommand{\stab}{\mathfrak{stab}}
\title{A stabilizer interpretation of the (extended) linearized double shuffle Lie algebra}
\author{Annika Burmester, Khalef Yaddaden}
\date{Source: arXiv:2603.05038v1 (CC BY 4.0)}
\begin{document}
\maketitle

\noindent\emph{Source and licence.} A stabilizer interpretation of the (extended) linearized double shuffle Lie algebra. Version arXiv:2603.05038v1 (CC BY 4.0), distributed by arXiv under the Creative Commons Attribution 4.0 International licence, http://creativecommons.org/licenses/by/4.0/, as recorded in the arXiv licence metadata for this exact version and shown on the abstract page https://arxiv.org/abs/2603.05038v1. Full licence notice: \texttt{licenses/arxiv-2603.05038-CC-BY-4.0.txt}.\par\smallskip

\noindent\textit{Editorial scope.} Counted unit: the article's theorem thm:stab\_equal\_lq (source0.tex lines 1085--1089). The article's complete original proof of it is delivered with it (source0.tex lines 1091--1144, one \texttt{proof} environment of 54 lines). No mathematics is written, repaired or re-derived: the statement and the proof are the article's own lines, in the article's own notation, and no retained line is substituted or re-flowed.  Retained uncounted as the source-local context the proof needs: lines 971--976; lines 1068--1073.  Nothing was omitted from any retained span, so the package carries no omission record.  No retained line contains a citation, so the wrapper carries no bibliography and no bibliography entry.  No retained line contains a figure environment, an image-inclusion command or drawing code, so no image is delivered and the package declares no asset dependency.  Editorial formatting only, in addition: an AMS \texttt{amsart} preamble that restates the article's own declarations and macros used by the retained lines, namely the environments \texttt{proposition}, \texttt{lemma}, \texttt{theorem}, \texttt{caselist} on separate counters, together with the standard packages those lines need.  The retained environments print in this document as Proposition 1, Lemma 1, Theorem 1, Caselist 1 in order of appearance, whereas the article prints the same items with the numbering of its own full document.  The complete native source of the article as distributed in the arXiv e-print for this exact version is bundled unchanged as \texttt{sources/195780/source0.tex}.
\begin{proposition} \label{prop:lq_subset_stab}
We have the following inclusion of subspaces of $\Lie(B)$
\[
\lqq\subset \stab_{\Lie(B)}(\tau).
\]
\end{proposition}

\begin{lemma} \label{lem:stab_subset_tau_inv}
We have the following equality of subspaces of $\Lie(B)$
\[
\stab_{\Lie(B)}(\tau)\subset\{\psi\in\Lie(B)\mid \tau(\pi_0(\psi))=\pi_0(\psi)\}.
\]
\end{lemma}

\begin{theorem} \label{thm:stab_equal_lq} We have the following equality of bigraded subspaces of $\Lie(B)$
\[
\stab_{\Lie(B)}(\tau)=\lqq\oplus \Q b_0.
\]
\end{theorem}

\begin{proof} The space $\Lie(B)$ is bigraded by weight and depth, where $b_0$ is of bidegree $(1,0)$ and $b_i$ is of bidegree $(i,1)$ for $i\geq1$. We consider several cases depending on the bidegree $(m,n)$.

\begin{caselist}
\item $(m,n)=(1,0)$. We show that
\begin{align} \label{eq:stab_equal_lq_(1,0)}
\stab_{\Lie(B)}(\tau)[1,0]=\lqq[1,0]\oplus\Q b_0.
\end{align}
By definition of $\lqq$, we have $\lqq[1,0]=0$. On the other hand, we have $\stab_{\Lie(B)}(\tau)[1,0]\subset \Lie(B)[1,0]=\Q b_0$. We compute $\sigma_{b_0}(w)=\ell_{b_0}w+\partial_{b_0}(w)=b_0w+[w,b_0]=wb_0$. Thus, the projection of $\sigma_{b_0}(w)$ to $\QB^0$ is always zero,
\[\sigma_{b_0}^0\equiv 0.\]
In particular, $\sigma_{b_0}^0$ commutes with $\tau$ and thus $b_0\in \Lie(B)$ is contained in $\stab_{\Lie(B)}(\tau)$. We obtain equality \eqref{eq:stab_equal_lq_(1,0)}. 
\item $m\geq 2$ and $n\geq2$. We have that
\[\stab_{\Lie(B)}(\tau)[m,n]=\lqq[m,n].\]
This follows from Proposition \ref{prop:lq_subset_stab} and Lemma \ref{lem:stab_subset_tau_inv} together with the observation that $\lqq$ and $\{\psi\in\Lie(B)\mid \tau(\pi_0(\psi))=\pi_0(\psi)\}$ coincide in bidegree $(m,n)$.
\item $m \geq 2$ is odd and $n=1$. By the same argumentation as in Case 2, we get
\[
\stab_{\Lie(B)}(\tau)[m,1]=\lqq[m,1].
\]
\item $m \geq 2$ is even and $n=1$. We show that
\[
\stab_{\Lie(B)}(\tau)[m,1]=\lqq[m,1].
\]
By definition of $\lqq$, we have $\lqq[m,1]=0$. On the other hand, by construction we have $\stab_{\Lie(B)}(\tau)[m,1]\subset \Lie(B)[m,1]=\bigoplus_{k=1}^m\Q \cdot \ad_{b_0}^{m-k}(b_k)$. Let $\psi=\sum_{k=1}^m \lambda_k \ad_{b_0}^{m-k}(b_k) $ for some $\lambda_1,\ldots,\lambda_m \in \Q$, and assume that $\psi\in \stab_{\Lie(B)}(\tau)$. We show that this implies $\psi=0$. Then, we have $\stab_{\Lie(B)}(\tau)[m,1]=0$ and hence the claimed equality.

\noindent
We compute 
\begin{align} \label{eq:sigma_circ_tau}
&\sigma_{\psi}^0\circ\tau(b_1)=\sigma_\psi^0(b_1) \\
&=\sum_{k=1}^m \lambda_k\pi_0\Big(\ad_{b_0}^{m-k}(b_k) b_1+\sum_{l=1}^k(-1)^{k+l}\binom{k-1}{l-1}[b_{1+k-l},\ad_{b_0}^{m-k}(b_l)]\Big) \nonumber\\
&=\sum_{k=1}^m \lambda_k \Bigg(\sum_{n=0}^{m-k} (-1)^{m-k-n} \binom{m-k}{n} b_0^nb_kb_0^{m-k-n}b_1+\sum_{l=1}^k (-1)^{k+l} \binom{k-1}{l-1}b_{1+k-l}b_0^{m-k}b_l \nonumber\\
&\hspace{0.4cm} - \sum_{k=1}^l\sum_{n=0}^{m-k}(-1)^{m+l-n} \binom{k-1}{l-1}\binom{m-k}{n}b_0^nb_lb_0^{m-k-n}b_{1+k-l}\Bigg), \nonumber
\end{align}
and thus
\begin{align} \label{eq:tau_circ_sigma}
\tau\circ\sigma^0_\psi(b_1)&=\sum_{k=1}^m \lambda_k \Bigg(\sum_{n=0}^{m-k} (-1)^{m-k-n} \binom{m-k}{n} b_{m-k-n+1}b_0^{k-1}b_{n+1}\\
&\hspace{0.4cm}+\sum_{l=1}^k (-1)^{k+l} \binom{k-1}{l-1}b_0^{l-1}b_{m-k+1}b_0^{k-l}b_1 \nonumber\\
&\hspace{0.4cm} - \sum_{k=1}^l\sum_{n=0}^{m-k}(-1)^{m+l-n} \binom{k-1}{l-1}\binom{m-k}{n}b_0^{k-l}b_{m-k-n+1}b_0^{l-1}b_{n+1}\Bigg). \nonumber
\end{align}
As we assumed $\psi\in\stab_{\Lie(B)}(\tau)$, the two expressions \eqref{eq:sigma_circ_tau} and \eqref{eq:tau_circ_sigma} must coincide. In particular, also the coefficient of any word
\[b_0^{t_1}b_sb_0^{t_2}b_1,\qquad s+t_1+t_2=m,\ t_1\geq1,\]
must agree in both expressions. In \eqref{eq:sigma_circ_tau} this coefficient is given by
\begin{align} \label{eq:coeff_sigma_circ_tau}
 \lambda_s \Big( (-1)^{t_2}\binom{t_1+t_2}{t_1}- (-1)^{t_2}\binom{t_1+t_2}{t_1}\Big)=0,  
\end{align}
and in \eqref{eq:tau_circ_sigma} this coefficient equals
\begin{align} \label{eq:coeff_tau_circ_sigma}
\lambda_{t_1+t_2+1}\Big((-1)^{t_2} \binom{t_1+t_2}{t_1}-(-1)^{m-t_2-1}\binom{t_1+t_2}{t_1}\Big)&=\Big(1-(-1)^{m-1}\Big) \lambda_{t_1+t_2+1}(-1)^{t_2} \binom{t_1+t_2}{t_1} \nonumber\\
&=2\lambda_{t_1+t_2+1}(-1)^{t_2} \binom{t_1+t_2}{t_1} 
\end{align}
The last equality follows from $m$ being even. As \eqref{eq:coeff_sigma_circ_tau} and \eqref{eq:coeff_tau_circ_sigma} must coincide, we deduce that $\lambda_{t_1+t_2+1}=0$. As $t_1\geq1,t_2\geq 0$ were arbitrarily chosen, we get
\[\lambda_k=0 \quad \text{ for } k=2,\ldots,m,\]
and hence $\psi=\lambda_1 \ad_{b_0}^{m-1}(b_1)$. For $m\geq2$, the element $\pi_0(\ad_{b_0}^{m-1}(b_1))=b_0^{m-1}b_1$ is not $\tau$-invariant, hence we deduce from Lemma \ref{lem:stab_subset_tau_inv} that $\lambda_1=0$. Thus, we get $\psi=0$.
\end{caselist}
As the bigraded spaces $\stab_{\Lie(B)}(\tau)$ and $\lqq\oplus\Q b_0$ coincide on all components, we get the desired equality.
\end{proof}

\end{document}
